\chapter{Arquitetura e Design}
\label{chap:arquitetura}

\section{Introdução}
\label{sec:arq-intro}

Este capítulo descreve as decisões arquiteturais e de design tomadas para o desenvolvimento da plataforma Gold Lock. São apresentados a arquitetura geral do sistema, o diagrama de componentes, o modelo de dados, o diagrama de classes e as decisões relativas ao design da interface.


\section{Visão Geral da Arquitetura}
\label{sec:arq-visao}

A plataforma segue uma arquitetura de microserviços, dividida em três camadas principais: camada de apresentação (frontend), camada de lógica de negócio (backend/\ac{API}) e camada de dados. A Figura~\ref{fig:arquitetura-componentes} apresenta o diagrama de componentes do sistema.

\begin{figure}[H]
\centering
\input{diagramas/componentes}
\caption{Diagrama de componentes da arquitetura do sistema.}
\label{fig:arquitetura-componentes}
\end{figure}


\section{Stack Tecnológica}
\label{sec:arq-stack}

A Tabela~\ref{tab:stack} apresenta as tecnologias selecionadas para cada camada do sistema, juntamente com a respetiva justificação.

\begin{table}[H]
\centering
\caption{Stack tecnológica do sistema Gold Lock.}
\label{tab:stack}
\small
\begin{tabularx}{\textwidth}{l l X}
\toprule
\textbf{Camada} & \textbf{Tecnologia} & \textbf{Justificação} \\
\midrule
Frontend Web & React 18 + TypeScript + TailwindCSS & Ecossistema maduro, tipagem forte, grande comunidade \\
Animações & Framer Motion & Necessário para efeitos Liquid Glass performantes \\
Backend & Node.js + Express + TypeScript & Stack unificada JS/TS, vasto ecossistema fintech \\
Base de Dados & PostgreSQL & Robusto, relacional, ideal para dados financeiros \\
Cache & Redis & Sessões, rate limiting, cache de dados Open Banking \\
IA/ML & Python (scikit-learn) & Classificação de transações e otimização fiscal própria, desenvolvida sem dependência de serviços externos \\
Open Banking & Salt Edge \ac{API} & Cobertura PT, sandbox gratuita, boa documentação \\
Autenticação & Supabase Auth & \ac{OAuth}2, \ac{2FA}, tier gratuito generoso \\
Containerização & Docker + Docker Compose & Ambiente reproduzível, isolamento de serviços, parity dev/prod \\
Deploy & Railway ou Render & Deploy simples, tier gratuito para \ac{MVP} \\
CI/CD & GitHub Actions & Automatização integrada com repositório \\
\bottomrule
\end{tabularx}
\end{table}


\section{Modelo de Dados}
\label{sec:arq-modelo-dados}

O modelo de dados foi desenhado em PostgreSQL, priorizando a integridade referencial e a normalização dos dados financeiros. A Figura~\ref{fig:modelo-er} apresenta o diagrama Entidade-Relacionamento do sistema.

\begin{figure}[H]
\centering
\input{diagramas/modelo-er}
\caption{Diagrama Entidade-Relacionamento do sistema Gold Lock.}
\label{fig:modelo-er}
\end{figure}

As principais entidades do modelo são:

\begin{itemize}
    \item \textbf{User:} Armazena informação do utilizador, incluindo credenciais e preferências.
    \item \textbf{BankAccount:} Representa uma conta bancária ligada via Open Banking, com referência à conexão Salt Edge.
    \item \textbf{Transaction:} Regista cada transação importada, incluindo descrição, montante, data e categoria.
    \item \textbf{Category:} Define as categorias de gastos, incluindo categorias pré-definidas para o mercado português.
    \item \textbf{Budget:} Armazena os orçamentos mensais definidos pelo utilizador por categoria.
    \item \textbf{SavingsGoal:} Representa as metas de poupança com valor alvo, prazo e progresso.
    \item \textbf{IRSSimulation:} Guarda os parâmetros e resultados das simulações de \ac{IRS}.
    \item \textbf{FiscalProfile:} Armazena os dados fiscais do utilizador introduzidos manualmente: rendimento bruto anual, contribuições para a Segurança Social, estado civil e número de dependentes. Estes dados são necessários para o cálculo personalizado do \ac{IRS} estimado.
    \item \textbf{DeductionAlert:} Regista os alertas gerados pelo sistema relativos a deduções identificadas nas transações (despesas dedutíveis detetadas), limites legais próximos e recomendações de otimização fiscal.
\end{itemize}


\section{Diagrama de Classes}
\label{sec:arq-classes}

A Figura~\ref{fig:diagrama-classes} apresenta o diagrama de classes simplificado do sistema, focando nas classes principais do domínio.

\begin{figure}[H]
\centering
\input{diagramas/classes}
\caption{Diagrama de classes do sistema Gold Lock.}
\label{fig:diagrama-classes}
\end{figure}


\section{Diagrama de Sequência}
\label{sec:arq-sequencia}

A Figura~\ref{fig:sequencia-openbanking} ilustra o diagrama de sequência para o fluxo de ligação de uma conta bancária via Open Banking, que constitui uma das interações mais críticas do sistema.

\begin{figure}[H]
\centering
\input{diagramas/sequencia-openbanking}
\caption{Diagrama de sequência: ligação de conta bancária via Open Banking.}
\label{fig:sequencia-openbanking}
\end{figure}


\section{Design da Interface --- Liquid Glass}
\label{sec:arq-design}

O design da interface segue o paradigma Liquid Glass, inspirado nas guidelines do iOS 26, adaptado ao contexto de uma aplicação financeira web. Os princípios fundamentais são:

\begin{enumerate}
    \item \textbf{Transparência e profundidade:} Utilização de efeitos de vidro fosco (\textit{backdrop-filter: blur}) em cards e painéis, criando sensação de profundidade e hierarquia visual.
    \item \textbf{Fluidez:} Animações suaves e responsivas implementadas com Framer Motion, priorizando 60 FPS em todas as interações.
    \item \textbf{Acessibilidade:} Cumprimento do nível AA das \ac{WCAG}, com rácio de contraste mínimo de 4.5:1 sobre efeitos de vidro.
    \item \textbf{Modo claro e escuro:} Suporte nativo a ambos os modos, com transição fluída entre eles.
\end{enumerate}

% TODO: Adicionar screenshots/mockups do Figma quando disponíveis


\section{Pipeline de IA de Otimização Fiscal}
\label{sec:arq-fiscal-ia}

O módulo de \ac{IA} de otimização fiscal é implementado como uma extensão do serviço \texttt{ml-service} existente, reutilizando a infraestrutura Python + Flask já disponível. Não é criado um container adicional --- toda a lógica corre no mesmo serviço Python, mantendo a arquitetura simples. O pipeline divide-se em três camadas funcionais independentes e sequenciais:

\begin{enumerate}
    \item \textbf{Classificador de Despesas Dedutíveis (ML):} Modelo Random Forest treinado com um conjunto de dados sintéticos derivados das regras do \ac{CIRS} 2025. Para cada transação da plataforma, o classificador determina se esta é dedutível para efeitos de \ac{IRS} e em que categoria (saúde, educação, habitação, encargos gerais familiares). Os dados de treino são gerados programaticamente a partir do mapeamento oficial entre tipos de comerciantes/despesas e os artigos 78.º-A a 78.º-F do \ac{CIRS}, garantindo a fidedignidade das labels.

    \item \textbf{Motor de Regras Fiscais CIRS 2025 (determinístico):} Implementação das regras fiscais portuguesas vigentes em Python puro, sem dependências externas. Inclui a tabela de escalões do \ac{OE} 2025 (Lei n.º~24-D/2024), os limites de dedução por categoria e a dedução das contribuições para a Segurança Social ao rendimento bruto (art.º 25.º \ac{CIRS}). O cálculo segue o fluxo: rendimento bruto $\rightarrow$ deduções ao rendimento (SS) $\rightarrow$ rendimento coletável $\rightarrow$ aplicação dos escalões $\rightarrow$ coleta bruta $\rightarrow$ deduções à coleta $\rightarrow$ imposto final. Os dados da \ac{SS} são introduzidos manualmente pelo utilizador no perfil fiscal (\texttt{FiscalProfile}), uma vez que a \ac{API} da Segurança Social Direta não é de acesso público.

    \item \textbf{Otimizador Fiscal (algoritmo de otimização):} Algoritmo próprio, implementado sem bibliotecas de otimização externas, que gera e compara múltiplos cenários de declaração possíveis: tributação individual vs.\ conjunta (para casais), inclusão ou exclusão de diferentes conjuntos de deduções, e ordenação das despesas pelas categorias com maior impacto fiscal. O algoritmo aplica uma estratégia \textit{greedy} com verificação exaustiva dos cenários de tributação conjunta/separada, devolvendo um ranking de cenários ordenado pelo imposto a pagar e uma lista de alertas sobre deduções não utilizadas ou próximas do limite legal.
\end{enumerate}

O fluxo de dados do pipeline é o seguinte:

\[
\text{Transações} \xrightarrow{\text{Classificador ML}} \text{Deduções identificadas} \xrightarrow{\text{Motor CIRS}} \text{IRS estimado} \xrightarrow{\text{Otimizador}} \text{Recomendações}
\]

Os novos endpoints expostos pelo \texttt{ml-service} e consumidos pelo backend Node.js são: \texttt{POST /fiscal/classify} (classificação de transações), \texttt{POST /fiscal/simulate} (cálculo do \ac{IRS} estimado) e \texttt{POST /fiscal/optimize} (comparação de cenários e recomendações).


\section{Conclusões}
\label{sec:arq-conclusoes}

A arquitetura proposta garante a separação clara de responsabilidades, a escalabilidade do sistema e a facilidade de manutenção. A escolha de uma stack unificada em TypeScript (frontend e backend) reduz a complexidade do desenvolvimento, enquanto a utilização de Python para o módulo de \ac{ML} e de otimização fiscal permite tirar partido do ecossistema científico da linguagem. O pipeline de \ac{IA} fiscal, inteiramente desenvolvido no projeto sem dependências de \ac{LLM}s externos, constitui o contributo técnico mais relevante da arquitetura, combinando aprendizagem automática supervisionada, motor de regras determinístico e algoritmo de otimização num fluxo coeso e auditável. O design Liquid Glass diferencia a plataforma visualmente dos concorrentes, contribuindo para uma experiência de utilizador premium.
